% easychair.tex,v 3.5 2017/03/15

\documentclass[a4paper]{easychair}

\usepackage{doc}
\usepackage{todonotes}

% In order to save space or manage large tables or figures in a
% landcape-like text, you can use the rotating and pdflscape
% packages. Uncomment the desired from the below.
%
% \usepackage{rotating}
% \usepackage{pdflscape}

%% Front Matter
%%
% Regular title as in the article class.
%
\title{Groupoid Containers in Locally Univalent 2,1-Categories}

\author{
Michele De Pascalis\inst{1}
\and
Niccol{\`o} Veltri\inst{2}
\and
Tarmo Uustalu\inst{3}\inst{4}
}

\institute{
Tallinn University of Technology, Estonia
\email{michde@taltech.ee}
\and
Tallinn University of Technology, Estonia
\email{niccolo@cs.ioc.ee}
\and
Tallinn University of Technology, Estonia
\and
Reykjav\'ik University, Iceland
\email{tarmo@ru.is}
}

\authorrunning{De Pascalis, Veltri and Uustalu}

\titlerunning{Groupoid Containers in Locally Univalent 2,1-Categories}

\usepackage{stmaryrd} %for Dana Scott brackets
\usepackage{newunicodechar}

% \newunicodechar{≃}{\ensuremath{\simeq}}
% \newunicodechar{λ}{\ensuremath{\lambda}}

\newcommand{\cont}[2]{#1 \mathbin{\triangleleft} #2}
% \DeclareMathOperator{\map}{map}
% \DeclareMathOperator{\Mon}{Mon}
% \DeclareMathOperator{\Mnd}{Mnd}
\DeclareMathOperator{\Endo}{Endo}
% \DeclareMathOperator{\Poly}{Poly}
% \DeclareMathOperator{\Nat}{Nat}
% \DeclareMathOperator{\Fin}{Fin}
% \DeclareMathOperator{\Id}{Id}
% \DeclareMathOperator{\ICMS}{ICMS}
% \DeclareMathOperator{\isICMM}{isICMM}
\DeclareMathOperator{\id}{id}
% \DeclareMathOperator{\PSh}{PSh}
% \DeclareMathOperator{\coPSh}{coPSh}
% \DeclareMathOperator{\St}{St}
% \DeclareMathOperator{\refl}{refl}
% \DeclareMathOperator{\Wr}{Wr}

\DeclareMathOperator{\idl}{idl}
\DeclareMathOperator{\idr}{idr}
\DeclareMathOperator{\assoc}{assoc}

% \newcommand{\letin}[2]{\mathsf{let\:} #1 \mathsf{\:in\:} #2}
% \newcommand{\ISt}[1]{\St^I_{#1}}
\newcommand{\catname}[1]{\mathsf{#1}}
\newcommand{\Set}{\catname{Set}}
\newcommand{\Cont}{\catname{Cont}}
\newcommand{\GpdCont}{\catname{GpdCont}}
% \newcommand{\ICMSCat}{\catname{ICMSCat}} % There must be a better way
% \newcommand{\Cat}{\catname{Cat}}
% \newcommand{\SetI}{\Set^I}
% \newcommand{\IC}{\catname{IC}}
% \newcommand{\ICI}{\IC_I}
% \newcommand{\C}{\mathcal{C}}
% \newcommand{\D}{\mathcal{D}}
\newcommand{\ext}[1]{\llbracket{#1}\rrbracket}
% \newcommand{\x}{\times}
% \newcommand{\ox}{\mathbin{\otimes}}
% \newcommand{\I}{\mathsf{I}}
% \newcommand{\tosim}{\xrightarrow{\sim}}
% \newcommand{\toI}{\to^I}
% \newcommand{\upa}{\mathbin\uparrow}
% \newcommand{\ula}{\mathbin\nwarrow}
% \newcommand{\ura}{\mathbin\nearrow}
% \newcommand{\iupa}[1]{\mathbin{\uparrow_{#1}}}
% \newcommand{\iula}[1]{\mathbin{\nwarrow_{#1}}}
% \newcommand{\iura}[1]{\mathbin{\nearrow_{#1}}}
% \newcommand{\Peidx}{P\mathsf{e}_{\equiv}}
% \newcommand{\conste}{k_e}
% \newcommand{\Pibub}{\bullet^P}
% \newcommand{\eP}{\e^P}
% \newcommand{\smoosh}[1]{\overline{#1}}
% \newcommand{\sbar}{\overline{s}}
% \newcommand{\consts}{(\lambda \_ . s)}
% \newcommand{\inv}[1]{{#1}^{-1}}
% \newcommand{\mapi}{\inv\map}

% \DeclareMathOperator{\lop}{\mathsf{left}}
% \DeclareMathOperator{\rop}{\mathsf{right}}
% \DeclareMathOperator{\inl}{\mathsf{inl}}
% \DeclareMathOperator{\inr}{\mathsf{inr}}
% \DeclareMathOperator{\fst}{\mathsf{fst}}
% \DeclareMathOperator{\snd}{\mathsf{snd}}
% \newcommand{\eqdef}{\mathbin{\overset{\mathsf{def}}{=}}}
% \newcommand{\semi}{\mathbin{\pmb ;}}
% \newcommand{\pair}[2]{(#1, #2)}
% \newcommand{\triple}[3]{(#1 , #2 , #3)}
% \newcommand{\Unit}{\mathbb 1}

% \newcommand{\fcart}{\sigma}
% \newcommand{\fvert}{\pi}

% \newcommand{\Bool}{\mathsf{Bool}}
% \newcommand{\true}{\mathsf{true}}
% \newcommand{\false}{\mathsf{false}}
\newcommand{\hGpd}{\mathsf{hGpd}}
\newcommand{\Type}{\mathsf{Type}}

% \DeclareMathOperator{\data}{data\ }
% \DeclareMathOperator{\where}{\ where}
% \DeclareMathOperator{\var}{var}
% \DeclareMathOperator{\app}{app}
% \DeclareMathOperator{\lam}{lam}

% \newcommand{\mact}{\mathbin{\blacktriangleright}}
% \newcommand{\IWr}{\Wr^{\mact}}


\begin{document}

\maketitle

\vspace{1.2em}

% \setcounter{tocdepth}{2}
% {\small
% 	\tableofcontents}

%%%%%%%%%%%
%% INTRO %%
%%%%%%%%%%%

It is a well known result in the study of containers \cite{Abbott2005} that
their interpretation \(\ext{-}\) as mappings between types yields a full and
faithful functor from the category \(\Cont\) of containers and container
morphisms, to the category \(\Endo(\Set)\) of \(\Set\) endofunctors. Moreover,
it has been shown \cite{Uustalu2017} that this functor is strong monoidal when
the two categories are equipped with the monoidal structures of container
substitution on one side, and of functor composition on the other.

Stated in this way, these results apply to the interpretations of type theory
where types stand for sets. However, it is also well known by now that more
complex interpretations are possible and worthwhile. When attempting to
replicate the above reasoning in Homotopy Type Theory
% cite
(an interpretation in which types stand for homotopy types), one has to add the
constraint that shapes and positions are set-truncated, lest the full
faithfulness result fail to extend.

In this work we adopt a more foundational approach, exploring a higher
categorical setting in which groupoid-truncated containers shall be rightful
inhabitants, where we expect to (and indeed will) find analogous properties.
% pretty dramatic I know

\begin{section}{Locally Univalent 2,1-Categories}
 
 We deal with the following notion of (2,1)-category, adapted from
 \cite{Finster2021}:
 
 A locally univalent (2,1)-category is a wild category whose \(\hom\)-types are
 \(\hGpd\)s. Explicitly:
 
 \begin{align*}
	 C_0              & : \Type                                                                                   \\
	 C_1(X , Y)       & : \hGpd                                          & \forall & X , Y : C_0                  \\
	 \id_X            & : C_1(X , X)                                     & \forall & X : C_0                      \\
	 f \star g        & : C_1(X , Z)                                     & \forall & X, Y, Z : C_0,               \\
	                  &                                                  &         & f : C_1(X, Y), g : C_1(Y, Z) \\
	 \idl_f           & : \id_X \star f \equiv f                         & \forall & X , Y : C_0, f : C_1(X,Y)    \\
	 \idr_f           & : f \star \id_Y \equiv f                         & \forall & X , Y : C_0, f : C_1(X,Y)    \\
	 \assoc_{f, g, h} & : (f \star g) \star h \equiv f \star (g \star h) & \forall & X
	 , Y , Z , W : C_0,                                                                                           \\
	 \textsf{etc.}
 \end{align*}
 \todo[inline]{put in the rest}
 
 
\end{section}

\begin{section}{Groupoid Containers}
 
\end{section}

\begin{section}{Monoids in \(\GpdCont\)}
 
\end{section}

\begin{section}{Free monoids in \(\GpdCont\)}
 
\end{section}


\label{sect:bib}

\bibliographystyle{plain}
\bibliography{main}

%------------------------------------------------------------------------------
\end{document}

